\scanpage{125}
% archive sheet number: 124
% source lines 1--11
semi-simples, on a une catégorie qui n'est plus graduée, seulement filtrée, \deleted{\unsure{I}} Si $V$ est la catégorie des schémas de type fini sur un corps fixé, $M$ la catégorie des motifs sur ce même corps (\unsure{présumée|pvrésumée} avec toutes les données et conditions de 1º)\handwritten{\deleted{\unsure{(}}}, on doit avoir un foncteur \deleted{\unsure{$V$|$Y$}} compatible avec ces structures \deleted{de} (cohomologie motivique)
\[
h:V\handwritten{\longrightarrow}\unsure{\mathbb{G}|G|E}\mathrm{rad}(M)\handwritten{;}\quad\text{(structure de Koko sur }\mathrm{Grad}(M)\text{)},
\]
et on désire que les foncteurs cohomologiques habituels sur $V$ se factorisant par \deleted{$M$. \unsure{On impose}}\inserted{des foncteurs sur $M$. On voudra} que les foncteurs \deleted{fibres} obtenus sur $M$ \unsure{soinnt|sointnt|soient} des foncteurs fibres \deleted{\uline{tordus} seulement} véritables\inserted{\handwritten{(non tordus)}}, ce qui n'est pas idiot, \deleted{car}\inserted{\handwritten{car}} on avait mis la structure de symétrie de Koko sur $\mathrm{Grad}(M)$: on s'est débarassé ainsi une fois pour toutes de la ``torsion''.

% source lines 12--15
On peut aussi, pour parler de vrais foncteurs fibres sur la catégorie des motifs semi-simples, remplacer la donnée de symmétrie ``géométrique'' i.e. de Koko pour le produit tensoriel dans \deleted{$\mathrm{Grad}(M)$}\inserted{la catégorie graduée $M$}, par la même canulée par le signe $(-1)^{ij}$ quand on a des \unsure{arguments} de degré $i,j$.

% source lines 16--24
11) Bien sûr, il y a lieu de chercher des théorèmes de structure intrinsèques, en termes de gerbes, pour des catégories $\handwritten{\unsure{\mathscr{C}|\mathfrak{C}|\mathbb{C}}}$ munies des structures de 1º,\inserted{avec $k=\mathrm{End}(\underline{1})$ un corps donné.} Il y a lieu au besoin de supposer qu'il existe une extension de $k$ sur laquelle il y ait un foncteur fibre \unsure{tordu|toodu} gradué (ou mltigradué ? Mais où s'arrêter ? ). Le plus beau serait évidemment de trouver un théorème de structure \unsure{ggnéral|ggénéral|général}, sans \uline{supposer} l'existence d'un foncteur fibre plus ou moins tordu où que ce soit, en définissant une notion assez générale de foncteur fibre tordu (?) et en ppuvant qu'\unsure{ip|il} y en a toujours sur une extension convenable de $k$.
